\section{Step 31: Call the record validator on scratch}
Form the record address from the file base plus the current offset. Pass
the saved record byte count, admitted scratch address, and private record-context
address to \texttt{video\_apply\_record}. This is a nested call, so arguments
and temporary registers may be overwritten. The outer cursor and frame count
stay in saved registers; the prologue has already saved this function's own
return address.

Fill \texttt{index\_validate}. Zero means the record passed its kind, sequence,
period, CRC and payload checks and was applied to scratch. Propagate development
status four if the prerequisite is unfinished. Any completed rejection takes the
index failure path. None of these paths publishes a directory entry yet.

\paragraph{Observe} Prepare \texttt{s31-1} and stop before
\texttt{index\_record\_valid}. At the first \texttt{index\_record\_valid}, the nested return
is zero and 16\,384. The private expected sequence at \texttt{\$sp+136} is
one. Scratch contains frame zero; the public directory and context are unchanged.

\Needspace{7\baselineskip}
\paragraph{Check} Which return address must the outer prologue preserve?
\begin{itemize}
  \item[A.] The checksum value.
  \item[B.] The next frame index.
  \item[C.] Its own caller's return address.
\end{itemize}
